\section{Обзор литературы}

Сходимость стохастических методов первого порядка подробно изучена~\citep{schneider2007stochastic, heyman2004stochastic, spall1998implementation, stich2019unified}, в то время как методы с предобуславливанием остаются относительно новыми и малоизученными. Идея адаптивного выбора шага восходит к онлайн-оптимизации~\citep{hazan2007adaptive}; в одной из первых работ по покоординатному предобуславливанию~\citep{duchi2011adaptive} авторы провели тщательный анализ теории сходимости Adagrad.
Однако в более поздних работах, например, обсуждающих RMSProp~\citep{rmsprop} или Adam~\citep{kingma2014adam}, теоретическим аспектам уделяется мало внимания, либо существующая теория содержит неточности в доказательствах.

Со временем ошибки были исправлены, что привело к разработке надёжных теорий сходимости для методов с предобуславливанием~\citep{reddi2019convergence, defossez2020simple}.
В другой работе Лощилов и Хуттер~\citep{loshchilov2017decoupled} изучили свойства алгоритмов Adam и AdamW в терминах гиперпараметров, а также рассмотрели методы рестартов. Чжан и др.~\citep{zhang2019lookahead} предложили оптимизатор Lookahead --- обёртку над произвольным базовым методом (в том числе Adam), в которой <<медленные>> веса обновляются по нескольким шагам <<быстрых>>. В работе~\citep{ginsburg2020stochastic} исследователи из Nvidia предложили оптимизатор NovoGrad с послойной адаптацией моментов и отделённым затуханием весов, показавший в их экспериментах прирост качества обучения.
Сравнительно недавно были построены теории сходимости для современных методов с предобуславливанием, таких как OASIS~\citep{jahani2021doubly, sadiev2022stochastic}.
Кроме того, появилась теория, рассматривающая изменяющиеся во времени матрицы предобуславливания~\citep{beznosikov2022scaled}.
Тем не менее, многие вопросы в этой области остаются без ответа.
Некоторые из них сформулированы в конце введения и рассматриваются в настоящей работе.
